\subsection{THE REAL LINE AND REAL NUMBERS}

DEFINITION 2. Let R denote the topological group which is the completion of the additive group Q of the rational line. The elements of R are called real numbers; as a topological space, R is called the real line; as a topological group, R is called the additive group of the real line.

We shall always identify Q with the dense subgroup of R to which it is canonically isomorphic. With this convention, every rational number is a real number. Every real number which is not rational is said to be irrational; we have seen in Chapter II, § 3, no. 3 that such numbers exist (we shall show this in another way in § 3, no. 3 of this chapter; see also Exercise 2 to § 2); hence (Chapter III, § 2, no. 1) the set CQ of irrational numbers is dense in R.

We shall show that the order structure of Q can be extended to R in such a way that the extended ordering is still compatible with the additive group structure of R:

PROPOSITION 1. The relation \(y - x \in \overline{\mathbf{Q}}_+\) is an ordering on \(\mathbf{R}\) which makes \(\mathbf{R}\) into a linearly ordered set; is compatible with the additive group structure on \(\mathbf{R}\) , and induces the ordering \(x \leqslant y\) on \(\mathbf{Q}\) .

We begin by showing that the relations \(y - x \in \overline{\mathbf{Q}}_+\) and \(z - y \in \overline{\mathbf{Q}}_+\) imply \(z - x \in \overline{\mathbf{Q}}_+\) . Indeed, the function \(x + y\) is continuous on \(\mathbf{R} \times \mathbf{R}\) , and therefore by (8) we have \(\overline{\mathbf{Q}}_+ + \overline{\mathbf{Q}}_+ \subset \overline{\mathbf{Q}}_+\) (Chapter I, § 2, no. 1, Theorem 1). Next, we shall show that the relations \(y - x \in \overline{\mathbf{Q}}_+\) and \(x - y \in \overline{\mathbf{Q}}_+\) imply \(x = y\) ; this will establish that \(y - x \in \overline{\mathbf{Q}}_+\) is an ordering on \(\mathbf{R}\) . It is enough to show that \(\overline{\mathbf{Q}}_+ \cap (-\overline{\mathbf{Q}}_+) = \{0\}\) . Now the functions \(x \to x^+\) and \(x \to x^-\) are uniformly continuous on \(\mathbf{Q}\) and can therefore be extended by continuity to \(\mathbf{R}\) (Chapter II, § 3, no. 6, Theorem 2); let \(f\) and \(g\) be their respective extensions. By extension we have \(x = f(x) - g(x)\) for all \(x \in \mathbf{R}\) ; if \(x \in \overline{\mathbf{Q}}_+\) then \(g(x) = 0\) , and since \(-\overline{\mathbf{Q}}_+\) is the closure of \(-\mathbf{Q}_+\) by the continuity of \(-x\) , it follows that if \(x \in -\overline{\mathbf{Q}}_+\) then \(f(x) = 0\) . Hence if \(x \in \overline{\mathbf{Q}}_+ \cap (-\overline{\mathbf{Q}}_+)\) we have \(f(x) = g(x) = 0\) and therefore \(x = 0\) .

By (10), we have \(\overline{\mathbf{Q}}_{+} \cup (-\overline{\mathbf{Q}}_{+}) = \mathbf{R}\) , and hence \(\mathbf{R}\) is linearly ordered by the ordering \(y - x \in \overline{\mathbf{Q}}_{+}\) .

Furthermore, since the relations \(y - x \in \overline{\mathbf{Q}}_+\) and \((y + z) - (x + z) \in \overline{\mathbf{Q}}_+\) are equivalent, the ordering \(y - x \in \overline{\mathbf{Q}}_+\) is compatible with the additive group structure of \(\mathbb{R}\) .

Finally, if \(x\) and \(y\) belong to \(\mathbf{Q}\) the relations \(y - x \in \overline{\mathbf{Q}}_+\) and \(y - x \in \mathbf{Q}_+\) are equivalent, and therefore the relation \(y - x \in \overline{\mathbf{Q}}_+\) induces the relation \(x \leqslant y\) on \(\mathbf{Q}\) . This completes the proof.

The relation \(y - x \in \overline{\mathbf{Q}}_+\) is again denoted by \(x \leqslant y\) . The set \(\overline{\mathbf{Q}}_+\) is the set of all \(x \geqslant 0\) in \(\mathbf{R}\) and is denoted by \(\mathbf{R}_+\) ; it is a closed set. The set of all \(x > 0\) is denoted by \(\mathbf{R}_+^*\) ; it is the complement of \(-\mathbf{R}_+\) and is therefore open in \(\mathbf{R}\) .

